\subsubsection*{Statement B.2t (Proved): abelian realization through a Galois alteration}

\begin{lemma}[Equivariant alteration data]
\label{B:equivariant-constant-data}
Let $k$ have characteristic $p>0$, let $K/k$ be regular and
finitely generated, and let $X_0$ be a proper geometrically
normal model.  We may choose a finite field extension $l/k$,
a smooth projective geometrically integral $l$-scheme $Y$, a
proper generically finite $k$-morphism $Y\to X_0$, and a finite
group $\Gamma$ acting on this morphism, with the following
properties.  Writing $L=k(Y)$, the field $l$ is stable under
$\Gamma$, its action on $Y$ is semilinear over this action on
$l$, and
\[
 K'=L^\Gamma\quad\hbox{is finite purely inseparable over }K.
\]
In particular $l_0=l^\Gamma$ is finite purely inseparable over
$k$, and $l/l_0$ is finite Galois.  The group action on $l$
need not be trivial.
\end{lemma}
\begin{proof}
Apply the regular Galois alteration theorem over $k^{\mathrm{perf}}$
to $X_0\times_k k^{\mathrm{perf}}$; use an initial projective
birational modification if needed.  The precise primary
statements are Situation~5.3, Theorem~5.13 and Corollary~5.15 of
\cite[pp.~613 and 619--620]{deJong1997}; projectivity is retained
in Theorem~5.9 and Proposition~5.11.  We obtain a regular
projective integral $\widetilde Y$ and a finite group $\Gamma$
with invariant function field purely inseparable over
$Kk^{\mathrm{perf}}$.  As the base field is perfect,
$\widetilde Y$ is smooth.  Let $\widetilde l$ be its field of
algebraic constants.  It is a finite separable extension of
$k^{\mathrm{perf}}$, and $\widetilde Y$ is geometrically
integral over $\widetilde l$.

The group preserves $\widetilde l$ and
$\widetilde l^\Gamma=k^{\mathrm{perf}}$.  Indeed an algebraic
constant in the invariant field is purely inseparable over
$Kk^{\mathrm{perf}}$ and separable over the perfect constant
field.  Regularity of $Kk^{\mathrm{perf}}/k^{\mathrm{perf}}$
forces this constant to lie in $k^{\mathrm{perf}}$.

Choose a finite extension $l/k$ inside $\widetilde l$ containing
a primitive generator of $\widetilde l/k^{\mathrm{perf}}$ and
all coefficients needed for the scheme, the map, the group
action and its finitely many action identities.  Adjoin their
finitely many $\Gamma$-conjugates.  Enlarging this finite field
if necessary, finite-presentation descent gives all the
specified data over $l$, and
$\widetilde l=lk^{\mathrm{perf}}$.  Smoothness, projectivity and
geometric integrality of $Y/l$ follow by faithfully flat field
descent.  Its map to $X_0$ is dominant of the same dimension,
hence generically finite, and is proper.  Thus $L/K$ is finite.

There is a compatible embedding $L\subset k(\widetilde Y)$.
Consequently $L^\Gamma$ lies in the purely inseparable extension
$k(\widetilde Y)^\Gamma$ of $Kk^{\mathrm{perf}}$.
Since $Kk^{\mathrm{perf}}/K$ is purely inseparable,
$L^\Gamma/K$ is purely inseparable; it is finite since $L/K$
is finite.  The same argument gives
$l^\Gamma=l\cap k^{\mathrm{perf}}$.  The final Galois assertion
is the fixed-field theorem for a finite group of field
automorphisms.  This construction expressly allows the finite
separable constants that an alteration over a perfect field
can introduce.
\end{proof}

\begin{lemma}[Realization of invariant abelian subgroups]
\label{B:invariant-abelian-realization}
Choose the preceding data and put
\[
 P=\operatorname{Pic}^{\tau}_{Y/l},\qquad
 \mathcal R=\operatorname{Res}_{l/k}(P),\qquad
 \mathcal A=\mathcal R^\Gamma.
\]
Let $H\hookrightarrow\mathcal A$ be an abelian variety over
$k$.  There is a torsor $Q$ over $k$ under an abelian variety,
and a rational map $X_0\dashrightarrow Q$, such that its
pullback on connected Picard schemes, evaluated on the fixed
alteration, has scheme-theoretic image exactly $H$ in
$\mathcal R$.

More precisely, let $\Gamma_0$ be the kernel of the action on
$l$, put $n=|\Gamma_0|$, and choose $q=p^e$ with $K'^q\subset K$.
We may take $Q$ to be a torsor under $(H^\vee)^{(q)}$, and the
resulting homomorphism is
\[
 H^{(q)}\xrightarrow{V_H^e}H\xrightarrow{[n]}H
       \hookrightarrow\mathcal R.                 \tag{B.2.14}
\]
Here $V_H^e$ is the dual of the $e$-fold relative Frobenius of
$H^\vee$; it is an isogeny.  In particular the image assertion
does not require either $n$ or $q$ to be prime to $p$.
\end{lemma}
\begin{proof}
\paragraph{Evaluation has finite kernel on $H$.}
Adjunction gives $u:H_l\to P$.  We check that its kernel is
finite.  Over an algebraic closure $\Omega$ of $k$, write
\[
 l\otimes_k\Omega=\prod_{\sigma}A_\sigma,
\]
where each $A_\sigma$ is an Artin local $\Omega$-algebra with
residue field $\Omega$.  Reduction defines a map from
$\mathcal R_\Omega$ to the product of the corresponding
geometric Picard schemes.  Its kernel is affine.  To see this
without a smoothness assertion about $P$, choose an affine
neighbourhood of the identity in each base-changed Picard
scheme over $A_\sigma$.  Every morphism from $T_{A_\sigma}$
whose reduction is the identity lands in that neighbourhood,
because $T\to T_{A_\sigma}$ is a nilpotent closed immersion
and a homeomorphism.  The kernel is therefore a closed
subscheme of the Weil restriction of an affine scheme along
a finite locally free map, which is affine.

The group $\Gamma$ acts transitively on the embeddings indexed
by $\sigma$, since $l^\Gamma/k$ is purely inseparable.
As $H$ lies in the pointwise fixed subgroup, the kernels of all
its geometric evaluation maps are equal.  Their common kernel
is contained in the affine reduction kernel just described.
It is also a closed subgroup of the proper scheme $H_\Omega$.
It is therefore proper and affine, and hence finite.
This proves the assertion for $u$ by field descent.

\paragraph{The Albanese torsor and its semilinear action.}
The image of $u$ is an abelian subvariety of $P$.
Duality for the proper geometrically integral scheme $Y/l$
gives a surjective homomorphism
$\operatorname{Alb}_{Y/l}\to H_l^\vee$.
Push out the Albanese torsor along this homomorphism.  We get
an $H_l^\vee$-torsor $T/l$ and a morphism $f:Y\to T$ whose
pullback on connected Picard schemes is $u$.
We use the proper geometrically reduced Albanese duality in
\cite[Theorem~3.2 and Remark~3.3, author pp.~12--13]{C_ACMV2026}.

The invariance $H\subset\mathcal R^\Gamma$ and functoriality
of the Albanese torsor give a semilinear $\Gamma$-action on
$T$ making $f$ equivariant.  Its linear action on the acting
abelian variety is the natural semilinear action on
$H^\vee_l$, because $\Gamma$ fixes the $k$-group $H$ pointwise.
In particular every $\gamma\in\Gamma_0$ acts on $T$ by a
translation $c_\gamma\in H^\vee(l)$.  The assignment
$\gamma\mapsto c_\gamma$ is a homomorphism: $\Gamma_0$ acts
trivially on $l$ and on $H^\vee$.  Thus $[n]c_\gamma=0$.

Push $T$ out along $[n]:H_l^\vee\to H_l^\vee$, obtaining
$T_n$ and $f_n:Y\to T_n$.  The subgroup $\Gamma_0$ acts
trivially on $T_n$.  The remaining semilinear action is the
Galois descent datum for $l/l_0$, where $l_0=l^\Gamma$.
It descends $T_n$ to an $H^\vee_{l_0}$-torsor $T_0/l_0$.
On function fields the equivariant map $f_n$ gives a
$\Gamma$-fixed point of $T_0(L)$, hence a point of $T_0(K')$.

\paragraph{Purely inseparable descent by a field homomorphism.}
The inclusion $l_0\subset K'$ and $K'^q\subset K$ imply
$l_0^q\subset k$: the field $k$ is algebraically closed in
$K$.  Hence $a\mapsto a^q$ is a field homomorphism
$l_0\to k$.  Define
\[
 Q=T_0\times_{\operatorname{Spec}l_0, a\mapsto a^q}
                   \operatorname{Spec}k.
\]
It is a torsor under $(H^\vee)^{(q)}$.
Base-changing the $K'$-point of $T_0$ along the field
homomorphism $K'\to K$, $b\mapsto b^q$, gives a $K$-point
of $Q$, and therefore a rational map $X_0\dashrightarrow Q$.

The canonical $k$-morphism $T_0\to Q$ has, on coordinate
algebras, the formula $b\otimes c\mapsto b^q c$.
After extension to $l_0$ it is the relative Frobenius on the
torsor.  Its pullback on connected Picard schemes is
$V_H^e:H^{(q)}\to H$.  Combining it with the pushout $[n]$
and the original adjunction homomorphism $u$ gives precisely
(B.2.14).  Both $V_H^e$ and $[n]$ are finite faithfully flat
isogenies.  Their composite has image all of $H$, proving
the claimed scheme-theoretic image in $\mathcal R$.
\end{proof}

\begin{lemma}[Relative Picard sheaves with finite constants]
\label{B:Pic-finite-constants-sheaf}
Let $l/k$ be finite and let $Z/l$ be proper and geometrically
integral, with a nonempty smooth open.  There is a canonical
isomorphism of full fppf Picard sheaves
\[
 \operatorname{Pic}_{Z/k}\simeq
 \operatorname{Res}_{l/k}\operatorname{Pic}_{Z/l}.
\]
No identification of their torsion components is asserted.
\end{lemma}
\begin{proof}
Choose a nonempty affine smooth open $U\subset Z$.
The affine scheme $S=\operatorname{Res}_{l/k}U$ is smooth
and surjective over $k$.  Indeed finite locally free restriction
preserves formal smoothness and finite presentation; over an
algebraic closure the algebra $l\otimes_k\bar k$ is a product
of Artin local algebras, and a point of each smooth nonempty
geometric fibre lifts through its nilpotent thickening.
The universal point gives a section of $Z_{S_l}\to S_l$.

After the fppf base change $T\times_kS\to T$, every class in
$\operatorname{Pic}_{Z/l}(T_l)$ is represented by a line bundle
on $Z_{T_l}$.  For clarity, the section permits rigidification:
the universal equality $H^0(Z,\mathcal O_Z)=l$ makes the
automorphism of a rigidified line bundle trivial, so the local
representatives and their rigidifications descend effectively.
Two representatives differ by a line bundle from $T_l$.
Such a line bundle becomes trivial after an fppf cover of $T$,
by the smooth-surjective generator-open argument of B.2m.
Thus sending a representative to its $k$-relative Picard class
is well-defined as a map of fppf sheaves and is inverse to the
natural map in the other direction.
\end{proof}

\begin{remark}[Application to the fixed comparison group]
Every comparison map $h_X$ of B.2m lands in
$\mathcal R^\Gamma$: the normal main graph $Y_X$ carries the
group action, its comparison with the regular $Y$ is
equivariant, and its pullback from $X$ is invariant.
Moreover the preceding full-sheaf identification applies to
every normal main alteration graph $Y_X/l$: its function field
is regular over $l$, so it has a nonempty smooth open.  Combining
it with B.1g over $l$ maps the entire comparison group
$\mathcal R$ into $\operatorname{Pic}_{Y_X/k}$, making the
generic-fibre restriction kernels used in the stabilization
argument defined on the full ambient group.
The construction retains the possible separable constants of an
alteration over $k^{\mathrm{perf}}$.  Its semilinear action,
translation quotient and purely inseparable field homomorphism
are part of the proof, rather than an assumption that those
constants disappear.  All three lemmas have undergone the
targeted proof checks recorded in the accompanying ledger.
\end{remark}
